\documentclass[12pt]{article}

\usepackage[a4paper, margin=1in]{geometry}
\usepackage{array}
\usepackage{longtable}
\usepackage{titlesec}
\usepackage{xcolor}
\usepackage{graphicx}
\usepackage{float}
\usepackage{hyperref}

% Define custom colours
\definecolor{headerblue}{RGB}{30,60,120}
\definecolor{veryminor}{RGB}{0,130,0}
\definecolor{minor}{RGB}{90,140,0}
\definecolor{significant}{RGB}{180,140,0}
\definecolor{major}{RGB}{200,100,0}
\definecolor{severe}{RGB}{180,0,0}
\definecolor{rare}{RGB}{0,130,0}
\definecolor{unlikely}{RGB}{90,140,0}
\definecolor{moderate}{RGB}{180,140,0}
\definecolor{high}{RGB}{200,100,0}
\definecolor{certain}{RGB}{180,0,0}

% Section styling
\titleformat{\section}
  {\large\bfseries\color{headerblue}}
  {}{0em}{}
\setlength{\parindent}{0pt}

\title{Authentication Endpoint: User Enumeration and Missing Rate Limiting}
\author{Meg Maroni}
\date{\today}

\begin{document}

\maketitle

\section*{Assessment Details}

\begin{tabular}{|p{5cm}|p{7.5cm}|}
\hline
Name & Meg Maroni \\ \hline
Email & \href{mailto:s221193238@deakin.edu.au}{s221193238@deakin.edu.au} \\ \hline
Team & AppAttack PenTesting \\ \hline
Role & Junior Lead \\ \hline
Project & EVAT \\ \hline
Quality Assurance &  Andrew Udawela\\ \hline
Was this Finding Successful? & Yes \\ \hline
Is this a Re-tested Finding? & No \\ \hline
\end{tabular}

\section*{Impact Values}

\begin{longtable}{|p{3cm}|p{12cm}|}
\hline
\textbf{\textcolor{veryminor}{Very Minor}} & Risk that holds little to no impact. Will not cause damage and regular activity can continue. \\ \hline
\textbf{\textcolor{minor}{Minor}} & \textcolor{minor}{Risk that holds minor impact. Can cause some damage but not enough to significantly affect operations.} \\ \hline
\textbf{\textcolor{significant}{Significant}} & Risk that causes noticeable disruption and may impede regular activity. \\ \hline
\textbf{\textcolor{major}{Major}} & Risk that significantly impacts operations and prevents the system from running normally. \\ \hline
\textbf{\textcolor{severe}{Severe}} & Risk that causes critical damage and may stop activity entirely. \\ \hline
\end{longtable}

\section*{Likelihood}

\begin{longtable}{|p{3cm}|p{12cm}|}
\hline
\textbf{\textcolor{rare}{Rare}} & Event may occur only in exceptional circumstances. \\ \hline
\textbf{\textcolor{unlikely}{Unlikely}} & Event could occur occasionally. \\ \hline
\textbf{\textcolor{moderate}{Moderate}} & \textcolor{moderate}{Event may occur under certain conditions.} \\ \hline
\textbf{\textcolor{high}{High}} & Event is likely to occur often. \\ \hline
\textbf{\textcolor{certain}{Certain}} & Event is occurring now or is expected to occur frequently. \\ \hline
\end{longtable}

\section*{Finding Details}

\subsection*{Finding Name: User Enumeration via Verbose Authentication Error Messages}

\subsubsection*{Description}
The login endpoint returns distinct, explicit error messages depending on whether a submitted email address corresponds to a registered account. When an existing account is submitted with an incorrect password, the API responds with a message confirming the account exists (e.g. \textit{"Authentication failed: Invalid password for email = [email.here]"}). When a non-existent email is submitted, the API instead responds with a message explicitly stating the account does not exist (e.g. \textit{"Authentication failed: The account with email = [email.here] does not exist"}). These two response bodies also differ in length (83 bytes vs. 118 bytes in testing), providing a second, independently detectable signal even if the message text were later obscured.

This finding is closely related to a separate finding submitted by another team member, included separately within this report and is referenced here only for context, which confirms that the same \texttt{/api/auth/login} endpoint applies no rate limiting, throttling, or account lockout even after 100 consecutive failed attempts. 

\subsubsection*{Risk}
An attacker in possession of a list of potential email addresses (e.g. sourced from public breach data, a company staff directory, or scraped contact lists) can submit each address to the login endpoint and, purely from the response message and/or response length, determine with certainty which addresses correspond to registered accounts on this platform, without needing to guess a password.

This enumeration finding compounds directly with the missing rate limiting reported, which allows the enumeration attempts to be run at a high volume without interruption, and the resulting list of confirmed valid accounts could then be fed into an unrestricted chain of password-guessing attempts against the same endpoint. Together, the two findings significantly reduce the effort required for an attacker to identify real users and subsequently compromise their accounts via brute-force or credential-stuffing attacks. The two findings should be considered jointly when assessing overall risk to this endpoint, as this increases the severity.

Independent of the rate-limiting issue, the enumeration signal itself also enables:
\begin{itemize}
    \item Targeted phishing campaigns against users known to hold an account on this platform.
\end{itemize}

\subsubsection*{Affected Assets}
\begin{itemize}
    \item EVAT WebApp: Authentication API (backend, Express/Node.js service)
    \item User account records exposed via the authentication response messages
\end{itemize}

\subsubsection*{Affected Path}
\texttt{POST /api/auth/login} (backend service, \texttt{http://localhost:8080})

\subsubsection*{Evidence}

\vspace{0.5em}

The following request/response pair was captured in Burp Suite when submitting a known, existing account email with an incorrect password:

\begin{verbatim}
HTTP/1.1 401 Unauthorized
Content-Type: application/json; charset=utf-8
Content-Length: 83

{"message":"Authentication failed: Invalid password for email = [Test1@gmail.com]"}
\end{verbatim}

\begin{figure}[H]
\caption{Utilising the known user account (Test1@gmail.com) with an incorrect password to monitor server response.}
\end{figure}
\begin{figure}[H]
\caption{Burp Suite Repeater response for a known, registered account. The server explicitly confirms the account exists.}
\end{figure}

\vspace{0.5em}

The same request was then repeated using an email address that does not correspond to any registered account:

\begin{verbatim}
HTTP/1.1 401 Unauthorized
Content-Type: application/json; charset=utf-8
Content-Length: 118

{"message":"Authentication failed: The account with email = 
[definately_not_a_real_email_1@gmail.com] does not exist"}
\end{verbatim}

\begin{figure}[H]
\caption{Utilising a non-registered account (definately\_not\_a\_real\_email\_1@gmail.com) to monitor server response.}
\end{figure}
\begin{figure}[H]
\caption{Burp Suite Repeater response for a non-existent account. The server explicitly confirms the account does not exist, and the response length differs from the known-account case (118 bytes vs. 83 bytes).}
\end{figure}

\vspace{0.5em}

\subsubsection*{Remediation Advice}
\begin{itemize}
    \item Return a single, generic error message for all authentication failures on both the UI and server response, regardless of whether the cause was an invalid email or an invalid password (e.g. \textit{"Invalid email or password"}), so the two cases cannot be distinguished.
    \item Ensure the response time and response body length are equivalent for both the "account does not exist" and "incorrect password" cases, to remove timing- and length-based side channels in addition to the message text.
    \item Apply the same generic-error principle to the registration endpoint and any password-reset/"forgot password" flow, as these are common alternative locations for the same class of enumeration vulnerability.
    \item Rate limiting and account lockout controls for this endpoint are addressed in the separate finding \textit{"Missing Rate Limiting on Login Endpoint Allows Unlimited Password Guessing"} (Shaharyar Nadeem); implementing both sets of remediations together is recommended, since either control alone leaves the other risk unmitigated.
\end{itemize}
\end{document}
